\section{USAMO 2026}
        \begin{enumerate}
        	\item (\textit{USAMO 2026}) Let $n$ be an integer greater than $1$. For which real numbers $x$ is 

 \[\lfloor nx\rfloor-\sum_{k=1}^n\frac{\lfloor kx\rfloor}{k}\]
maximal, and what is the maximal value that this expression can take?
Note: $\lfloor z\rfloor$ denotes the greatest integer less than or equal to $z$.


 



\textbf{Solution 1}

Let the given function be defined as:
 \[f(x) = \lfloor nx \rfloor - \sum_{k=1}^{n} \frac{\lfloor kx \rfloor}{k}\]

 First, we show that $f(x)$ is periodic with period 1. Let $x = y + 1$. Then:\begin{align*}f(x+1) &= \lfloor n(x+1) \rfloor - \sum_{k=1}^n \frac{\lfloor k(x+1) \rfloor}{k} \\&= \lfloor nx + n \rfloor - \sum_{k=1}^n \frac{\lfloor kx + k \rfloor}{k} \\&= \lfloor nx \rfloor + n - \sum_{k=1}^n \left( \frac{\lfloor kx \rfloor}{k} + 1 \right) \\&= \lfloor nx \rfloor + n - \sum_{k=1}^n \frac{\lfloor kx \rfloor}{k} - n \\&= f(x)\end{align*}
Because of this periodicity, we can restrict our analysis to $x \in [0, 1)$ and then generalize the result to all real numbers.

 We partition the interval $[0, 1)$ into $n$ subintervals of the form $I_m = \left[\frac{m}{n}, \frac{m+1}{n}\right)$ for $m = 0, 1, \dots, n-1$.

 Let $x \in I_m$. This implies $m \le nx < m+1$, so $\lfloor nx \rfloor = m$. Additionally, since $x \ge \frac{m}{n}$ and $k > 0$, we have $kx \ge \frac{km}{n}$, which means $\lfloor kx \rfloor \ge \left\lfloor \frac{km}{n} \right\rfloor$ for all $k \in \{1, \dots, n\}$.

 Since $\lfloor nx \rfloor$ remains constant at $m$ throughout $I_m$ and the terms $\lfloor kx \rfloor$ are being subtracted, the maximum of $f(x)$ on $I_m$ is bounded above by its value at the lower endpoint:
 \[f(x) = m - \sum_{k=1}^n \frac{\lfloor kx \rfloor}{k} \le m - \sum_{k=1}^n \frac{\left\lfloor \dfrac{km}{n} \right\rfloor}{k} = f\left(\dfrac{m}{n}\right)\]This proves that the global maximum of $f(x)$ on $[0, 1)$ must be attained at one of the points $x = \frac{m}{n}$.

 We now evaluate $f\left(\dfrac mn\right)$. Using the identity $\lfloor z \rfloor = z - \{z\}$ (where $\{z\}$ is the fractional part of $z$), we get:
 \[f\left(\dfrac{m}{n}\right) = m - \sum_{k=1}^n \dfrac{\dfrac{km}{n} - \left\{\dfrac{km}{n}\right\}}{k} = m - \sum_{k=1}^n \dfrac{m}{n} + \sum_{k=1}^n \frac{\left\{\dfrac{km}{n}\right\}}{k}\]Since $\sum_{k=1}^n \frac{m}{n} = m$ and $\left\{\frac{nm}{n}\right\} = 0$, this simplifies beautifully to:
 \[f\left(\dfrac{m}{n}\right) = \sum_{k=1}^{n-1} \frac{\left\{\dfrac{km}{n}\right\}}{k} = \frac{1}{n} \sum_{k=1}^{n-1} \frac{km \bmod n}{k}\]Let $c_k = km \bmod n$. We want to find the integer $m \in \{0, \dots, n-1\}$ that maximizes the sum $\sum_{k=1}^{n-1} \frac{c_k}{k}$.

 We claim that the sum is strictly maximized when $m = n-1$. For $m = n-1$, the sequence becomes $c_k = k(n-1) \bmod n = -k \bmod n = n-k$. 

 To prove that any other $m$ yields a smaller sum, we compare the sequence $c_k$ with the sequence $b_k = n-k$. Let $C$ be the multiset $\{c_1, \dots, c_{n-1}\}$ and $B$ be the set $\{1, 2, \dots, n-1\}$. 

 Let $d = \gcd(m, n)$. The multiset $C$ consists of the multiples of $d$ from $0$ to $n-d$. Specifically, $0$ appears $d-1$ times, and each multiple $vd$ (for $v \ge 1$) appears $d$ times. 

 For any integer $y \in \{1, \dots, n-1\}$, let $N_C(y)$ be the number of elements in $C$ that are $\ge y$. The elements $vd \ge y \implies v \ge \left\lceil \dfrac yd \right\rceil$. The number of such valid multipliers $v$ up to $\dfrac{n}{d}-1$ is exactly $\dfrac{n}{d} - \left\lceil \dfrac yd \right\rceil$. Since each non-zero multiple appears $d$ times, we have:
 \[N_C(y) = d \left(\frac{n}{d} - \left\lceil \dfrac yd \right\rceil\right) = n - d\left\lceil \dfrac yd \right\rceil\]Because $d\left\lceil \dfrac yd \right\rceil \ge d\left(\dfrac yd\right) = y$, we obtain $N_C(y) \le n - y$. The number of elements in $B$ that are $\ge y$ is precisely $N_B(y) = n - y$. Thus, $N_C(y) \le N_B(y)$ for all $y$.

 This establishes that the $i$-th largest element of $C$ is $\le$ the $i$-th largest element of $B$ (which is $n-i$). Consequently, for any $j \in \{1, \dots, n-1\}$, the sum of the first $j$ elements of the sequence $c_k$ is bounded by the sum of the $j$ largest elements of $C$, giving:
 \[\sum_{k=1}^j c_k \le \sum_{k=1}^j (n-k)\]

 Let $A_j = \sum_{k=1}^j (n-k - c_k)$. We have proven that $A_j \ge 0$ for all $j$. Using summation by parts on our sum:
 \[\sum_{k=1}^{n-1} \frac{n-k - c_k}{k} = \sum_{j=1}^{n-2} A_j \left(\frac{1}{j} - \frac{1}{j+1}\right) + A_{n-1} \frac{1}{n-1}\]Since $\frac{1}{j} - \frac{1}{j+1} > 0$ and $A_j \ge 0$, the entire sum is $\ge 0$. This implies $\sum_{k=1}^{n-1} \frac{c_k}{k} \le \sum_{k=1}^{n-1} \frac{n-k}{k}$. Equality holds if and only if $A_j = 0$ for all $j$, meaning $c_k = n-k$ for all $k$, which occurs only when $c_1 = m = n-1$.

 Since $m = n-1$ is the unique maximizer among our interval lower bounds, the maximum value of $f(x)$ evaluates to:
 \[f\left(\frac{n-1}{n}\right) = \frac{1}{n} \sum_{k=1}^{n-1} \frac{n-k}{k} = \sum_{k=1}^{n-1} \frac{1}{k} - \frac{n-1}{n} = \sum_{k=1}^n \frac{1}{k} - 1\]

 To find the continuous values of $x$ where this maximum is achieved, we inspect the interval $I_{n-1} = \left[\frac{n-1}{n}, 1\right)$. For any $x$ in this interval, and for $1 \le k \le n$:
 \[k - \frac{k}{n} = k \left(\frac{n-1}{n}\right) \le kx < k\]Because $1 \le k \le n$, the term $\frac{k}{n} \le 1$, meaning $k-1 \le k - \frac{k}{n}$. Therefore, $kx \in [k-1, k)$, which locks $\lfloor kx \rfloor$ firmly at $k-1$. Computing $f(x)$ anywhere in this interval yields the constant maximum:
 \[f(x) = (n-1) - \sum_{k=1}^n \frac{k-1}{k} = (n-1) - \sum_{k=1}^n \left(1 - \frac{1}{k}\right) = \sum_{k=1}^n \frac{1}{k} - 1\]

 \textbf{Conclusion:}

 Accounting for the periodicity of 1 established in Step 1, the maximal value is sustained across all intervals shifted by integers.

 
                \begin{itemize}
                    \item Maximal value:
                \end{itemize}


                
 \[\sum_{k=1}^n \frac{1}{k} - 1\]

 
                \begin{itemize}
                    \item Real numbers $x$ that achieve this: 
                \end{itemize}


                
 \[x \in \bigcup_{p \in \mathbb{Z}} \left[p - \frac{1}{n}, p\right)\]

 



\textbf{Solution 2}

After some manipulation (see Solution 1), let $x\in \left[\frac{m}{n}, \frac{m+1}{n}\right)$, where $m\in\{0,1,2,\ldots, n-1\}$. We proceed by maximizing
 \[f(x)\le f\left(\frac{m}{n}\right)=\frac{1}{n}\sum_{k=1}^{n}\frac{km\bmod{n}}{k}\]The maximal is clearly not achievable when $m=0$. Let $a_k=km\bmod n$.

 \textbf{Case 1: }$\gcd(m,n)=1$. 

 We note that:

 1) For $i,j\in\{0,1,\ldots,n-1\}, i\neq j$, we have $a_i\neq a_j \,$. To prove, if $a_i=a_j\,$ , then $a_i-a_j\equiv m(i-j)\equiv 0\pmod{n}$, which yields $i=j$ as a contradiction.

 2) The sequence $b_k=\frac{1}{k}$ is strictly decreasing for $k=1,2,\ldots, n$ as the summation occurs. By Rearrangement Inequality, the summation obtains its maximal when $a_1>a_2>\ldots > a_n$. 

 We have $a_k\in\{0,1,2,\ldots, n-1\}$. The sole way for those numbers to be strictly decreasing is
 \[n-1>n-2>\ldots>0\]In the case of maximal, this forces $a_k=n-k$.

 Thus,\begin{gather*}a_k\equiv km\equiv -k\pmod{n}\iff\\m\equiv -1\pmod{n}\end{gather*}We have $m=n-1$. This gives the maximal of $f$:
 \[\frac{1}{n}\sum_{k=1}^{n}\frac{km\bmod{n}}{k}\le \frac{1}{n}\sum_{k=1}^{n}\frac{n-k}{k}=\sum_{k=2}^{n}\frac{1}{k}\]

 \textbf{Case 2: }$\gcd(m,n)=d>1$

 Let $n=d\cdot q, m=d\cdot b$, where $\gcd{(q,b)}=1$. We get $a_k=km\bmod{n}=d\, (bk\bmod{q})$. This gives a multiset of values of $a_k$:
 \[a_k\in\bigl\{\underbrace{n-d, n-d,\ldots, n-d}_{d\text{'s }n-d}, \enspace\underbrace{n-2d, n-2d,\ldots, n-2d}_{d\text{'s }n-2d}\enspace ,\ldots,\enspace \underbrace{0, 0,\ldots, 0}_{d\text{'s }0}\bigr\}\]By Rearrangement Inequality, let $a_1^*\ge a_2^*\ge \ldots \ge a_n^*$ be a rearrangement of the sequence $a$ that yields the maximal. We have
 \[\frac{1}{n}\sum_{k=1}^{n}\frac{a_k}{k}\le\frac{1}{n}\sum_{k=1}^{n}\frac{a_k^*}{k}\]Generally, let $1\le t\le q, 1\le s\le d$. For indices $k=(t-1)d+s$, we have $a_k^*=n-td$. Note that
 \[a_k^*=n-td\le n-td+d-s=n-(t-1)d-s=n-k\]But the equality holds if and only if $s=d$. We see that $d\nmid k\implies a_k^* < n-k$, which means the upper bound from \textbf{Case 1} is \textbf{not attainable} for \textbf{Case 2}:
 \[\frac{1}{n}\sum_{k=1}^{n}\frac{a_k^*}{k} < \frac{1}{n}\sum_{k=1}^{n}\frac{n-k}{k}=\sum_{k=2}^{n}\frac{1}{k}\]

 Therefore, by our analysis from \textbf{Case 1}, the maximal is achieved when $m=n-1$:
 \[\frac{1}{n}\sum_{k=1}^{n}\frac{km\bmod{n}}{k}\le\sum_{k=2}^{n}\frac{1}{k}\]

 Furthermore, for $x=1-\frac{1}{n}+c$, where $0\le c < \frac{1}{n}$, since $-1\le k\left(c-\frac{1}{n}\right)<0$, we have\begin{align*}f(x)&=\lfloor nx\rfloor -\sum_{k=1}^{n}\frac{\lfloor kx\rfloor}{k}\\&=\lfloor n-1+cn\rfloor-\sum_{k=1}^{n}\frac{\left\lfloor k+k\left(c-\frac{1}{n}\right)\right\rfloor}{k}\\&=n-1-\sum_{k=1}^{n}\frac{k-1}{k}\\&=\sum_{k=2}^{n}\frac{1}{k}\\\end{align*}
We see that all real $x\in \left[1-\frac{1}{n}, 1\right)$ yields the maximal.

 $f(x)$ has a period of 1. Thus, the real $x$ that attains the maximal satisfies
 \[x\in\bigcup_{j\in\mathbb{Z}}\left[j-\frac{1}{n}, \, j\right)\]

 


	\item (\textit{USAMO 2026}) Annie is playing a game where she starts with a row of positive integers, written on a blackboard, each of which is a power of $2$. On each turn, she can erase two adjacent numbers and replace them with a power of $2$ that is greater than either of the erased numbers. This shortens the row of numbers, and she continues to take turns until only one number remains. Annie wins the game if the final remaining number is less than $4$ times the sum of the original numbers. Is it always possible for Annie to win, regardless of the starting row of numbers?


 



\textbf{Solution 1}

We start by realizing that the algorithm of combining the smallest term with its smallest neighbor seems to work, and we attempt to prove that this algorithm yields a sum which is $\le 4\cdot(\text{the original sum})$.

 We proceed with strong induction with the base cases of $n=1$ and $n=2$ (if $n=1$ our result is trivial and if $n=2$ if the two numbers are $x_1$ and $x_2$, the combined term will be $2\max(x_1,x_2)\le 2(x_1+x_2)$).

 Denote the original $n$ integers as $x_1,x_2,\dots,x_n$, and let
 \[S=2x_1+4x_2+4x_3+4x_4+\cdots+4x_{n-2}+4x_{n-1}+2x_n.\]

 We attempt to prove that our result is always $\le k$.

 Let $k$ be the first index such that
 \[2x_1+4x_2+\cdots+4x_{k-1}+4x_k=S_k\ge \frac{S}{2},\]and assume that $k\ne 1$ or $n$.

 WLOG, we can assume that $S_k-2x_k\le \frac{S}{2}$, as if not we can just mirror the sequence as both $S_k-2x_k$ and $S-S_{k-1}-2x_k$ cannot be greater than $\frac{S}{2}$, which contradicts our initial definition for $S$.

 By our inductive hypothesis, as
 \[S_k-2x_k=2x_1+4x_2+\cdots+4x_{k-1}+2x_k,\]we know we can combine the integers $x_1$ to $x_k$ to form an integer $\le \frac{S}{2}$, and similarly, as we know that the terms from $x_{k+1}$ to $x_n$ can be expressed an integer
 \[\le 2x_{k+1}+4x_{k+2}+\cdots+4x_{n-1}+2x_n\le S-S_k\le \frac{S}{2},\]implying that the combination of both terms $\le S$, as desired.

 In the case that $k=n$ (or if $k=1$ we just mirror the entire sequence), we consider the sum of the first $n-1$ terms, and as we know $S_{n-1}-2x_{n-1}\le \frac{S}{2}$, we can combine the first $n-1$ terms to a value which is $\le \frac{S}{2}$.

 As we know $x_n\le \frac{S}{2}$ (as $S\ge 2x_n$) we know that we can combine both terms to form a value $\le S$, as desired.

 



\textbf{Solution 2}

Let $X = (x_1, x_2, \dots, x_n)$ where $\forall i, x_i \in \{2^k \mid k \in \mathbb{Z}\}$, and let $S =\sum_{i=1}^n x_i$. 

 Define the adjacent merge operation $\omega(a, b) = 2\max(a, b)$. 

 Sequential adjacent applications of $\omega$ on $X \iff$ Construction of a full alphabetic (left-to-right) binary tree $T$.

 
                \begin{itemize}
                    \item $\forall \text{ leaf } i \in T$: $\text{value} = x_i$, $\text{depth} = d_i$.
                \end{itemize}


                
                \begin{itemize}
                    \item $\forall \text{ internal node with children } u, v$: $\text{value} = \omega(u,v) = \max(2u, 2v)$.
                \end{itemize}


                By induction, the root value $R$ is given by:
 \[R = \max_{1 \le i \le n} (x_i 2^{d_i})\]Winning condition $W \iff \exists \text{ valid tree } T \text{ such that } R < 4S$.

 $\because \forall i, x_i \in \{2^k\} \land \text{closure of } \omega \implies R \in \{2^k \mid k \in \mathbb{Z}\}$.

 Let $M = \max \{2^k \mid 2^k < 4S, k \in \mathbb{Z}\}$.

 \begin{align*}\because R \in \{2^k\} \implies W &\iff R \le M \\&\iff \forall i, \, x_i 2^{d_i} \le M \\&\iff \forall i, \, 2^{d_i} \le \frac{M}{x_i} \\&\iff \forall i, \, d_i \le \log_2 M - \log_2 x_i\end{align*}


 Let $L_i = \log_2 M - \log_2 x_i$. ($\because M, x_i \in \{2^k\} \implies L_i \in \mathbb{Z}$). 

 $\implies W \iff \exists \text{ alphabetic binary tree } T \text{ where } \forall i, d_i \le L_i$.

 Evaluate the Kraft sum constraint for $\{L_i\}$:
 \[\sum_{i=1}^n 2^{-L_i} = \sum_{i=1}^n 2^{-\log_2(M/x_i)} = \sum_{i=1}^n \frac{x_i}{M} = \frac{S}{M}\]By definition of $M$, we have $M < 4S \le 2M$.
 \[\implies \frac{1}{4} < \frac{S}{M} \le \frac{1}{2} \implies \sum_{i=1}^n 2^{-L_i} \le \frac{1}{2}\]

 Let $q_i = 2^{-L_i+1} = 2 \cdot 2^{-L_i} \implies\sum_{i=1}^n q_i = 2 \sum_{i=1}^n 2^{-L_i} \le 1$. 

 Define $C_1 = 0, C_{i+1} = C_i + q_i$, and intervals $I_i = [C_i, C_{i+1})$. 

 $\because\sum q_i \le 1 \implies \forall i, I_i \subseteq [0, 1] \land (i \neq j \implies I_i \cap I_j = \emptyset)$.

 Notice $|I_i| = 2 \cdot 2^{-L_i}$. By the fundamental property of dyadic intervals:
 \[\forall I \text{ with } |I|=2\cdot 2^{-L}, \exists \text{ sub-interval } J = [k \cdot 2^{-L}, (k+1) \cdot 2^{-L}) \subseteq I, \, k \in \mathbb{Z}\]Let $k_i = \lceil C_i \cdot 2^{L_i} \rceil$, and define $J_i = [k_i \cdot 2^{-L_i}, (k_i+1) \cdot 2^{-L_i})$.

 \begin{align*}&\implies J_i \subseteq I_i \\&\because I_i \cap I_j = \emptyset \land \text{strictly ordered left-to-right} \implies J_i \cap J_j = \emptyset \land \text{strictly ordered} \\&\because J_i = [k_i 2^{-L_i}, (k_i+1)2^{-L_i}) \iff \text{Unique binary prefix string of length } L_i \\&\implies \{J_i\}_{i=1}^n \text{ forms a valid alphabetic prefix code} \\&\implies \exists \text{ binary tree } T' \text{ with initial nodes } x_i \text{ at strict depth } L_i\end{align*}


 Let $T = \text{Prune}(T')$ (recursively eliminate $\text{degree-1}$ internal nodes, redirecting child to parent). 

 $\implies T \text{ is a full alphabetic binary tree}$. 

 $\because \text{Pruning strictly decreases or maintains depth without violating left-to-right order}$. 

 $\implies \forall i \in T, \text{ final depth } d_i \le L_i      \implies R = \max_{i=1}^n (x_i 2^{d_i}) \le \max_{i=1}^n (x_i 2^{L_i}) = M < 4S      \implies W \text{ is mathematically guaranteed}.$ $\square$

 


	\item (\textit{USAMO 2026}) Let $ABC$ be an acute scalene triangle with no angle equal to $60^\circ$. Let $\omega$ be the circumcircle of $ABC$. Let $\triangle_B$ be the equilateral triangle with three vertices on $\omega$, one of which is $B$. Let $\ell_B$ be the line through the two vertices of $\triangle_B$ other than $B$. Let $\triangle_C$ and $\ell_C$ be defined analogously. Let $Y$ be the intersection of $AC$ and $\ell_B$, and let $Z$ be the intersection of $AB$ and $\ell_C$.


 Let $N$ be the midpoint of minor arc $BC$ on $\omega$. Let $\mathcal{R}$ be the triangle formed by $\ell_B$, $\ell_C$, and the tangent to $\omega$ through $N$. Prove that the circumcircle of $AYZ$ and the incircle of $\mathcal{R}$ are tangent.


 



\textbf{Solution}

We establish a complex plane with $\omega$ being the unit circle, where $\lvert a\rvert=1,\lvert b\rvert=1,\lvert c\rvert=1$ represent points $A,B,C$, respectively. Let $B_1=e^{i\frac{2\pi}{3}}b$ and $B_2=e^{i\frac{4\pi}{3}}b$ be points from $\omega\cap\ell_B$. Also, note that we can assume $n=\sqrt{bc}$ representing $N$. Let $T_N$ be the tangent of $\omega$ through $N$. Since $B_1\cdot B_2=b^2$ and $B_1+B_2=-b$, we obtain the following equations:\begin{align*}\ell_B&: z+b^2 \overline{z}=-b\\\ell_C&: z+c^2 \overline{z}=-c\\AB&: z+ab\overline{z}=a+b\\AC&: z+ac\overline{z}=a+c\\T_N&: z+bc\overline{z}=2\sqrt{bc}\end{align*}
Solving for $D,E,F,Y,Z$, we have
 \[d=\frac{b(c+2n)}{b-c},\quad e=\frac{c(b+2n)}{c-b},\quad f=-\frac{bc}{b+c}\]
 \[y=\frac{b(ab+bc+ca)}{b^2-ac},\quad z=\frac{c(ab+bc+ca)}{c^2-ab}\]
 \[\overline{y}=-\frac{a+b+c}{b^2-ac},\quad\overline{z}=-\frac{a+b+c}{c^2-ab}\]We proceed by determining the incenter $I$ of $\mathcal{R}$. Note that
 \[\lvert d-f\rvert=\left\lvert \frac{2n(n^2+nc+c^2)}{b^2-c^2}\right\rvert =\left\lvert \frac{2n(n+b+c)}{b^2-c^2}\right\rvert\]By symmetry, we have $\lvert d-f\rvert=\lvert e-f\rvert$. Also, 
 \[\lvert d-e\rvert=\left\lvert \frac{2bc+2nc+2nb}{b-c}\right\rvert=\left\lvert \frac{2n(n+b+c)}{b-c}\right\rvert\]We see that the ratio of side $de$ to the other legs in the same isosceles triangle is $\lvert b+c\rvert=\sqrt{(b+c)\left(\overline{b}+\overline{c}\right)}=\frac{b+c}{n}$. Thus,
 \[I=\frac{d+e+\frac{b+c}{n}\cdot f}{1+1+\frac{b+c}{n}}=\frac{bc}{2n+b+c}\]Let $r$ be the inradius of the same circle. We have
 \[r=\left\lvert\frac{n^2}{2n+b+c}-n\right\rvert=\left\lvert\frac{n+b+c}{2n+b+c}\right\rvert\]We know that $r=\frac{n+b+c}{2n+b+c}$ because in that case, $r=\overline{r}\in\mathbb{R}$.

 Let $O$ be the circumcenter of $\triangle ayz$. We have
 \[\lvert O-a\rvert=\lvert O-y\rvert=\lvert O-z\rvert=R\]Squaring gives $O\overline{O}-a\overline{O}-\overline{a}O+1=R^2$. By symmetry, we get\begin{align*}(y-a)\overline{O}+(\overline{y}-\overline{a})O=y\overline{y}-1\\(z-a)\overline{O}+(\overline{z}-\overline{a})O=z\overline{z}-1\end{align*}
Note that
 \[b(a+b+c)(ab+bc+ca)+(b^2-ac)^2=(a^2+ab+b^2)(b^2+bc+c^2)\]It follows that\begin{align*}\frac{c(a^2+ab+b^2)}{b^2-ac}\overline{O}-\frac{a^2+ab+b^2}{a(b^2-ac)}O=-\frac{(a^2+ab+b^2)(b^2+bc+c^2)}{(b^2-ac)^2}\\\frac{b(a^2+ac+c^2)}{c^2-ab}\overline{O}-\frac{a^2+ac+c^2}{a(c^2-ab)}O=-\frac{(a^2+ac+c^2)(b^2+bc+c^2)}{(c^2-ab)^2}\\\end{align*}
Solving yields
 \[\overline{O}=-\frac{(a+b+c)(b^2+bc+c^2)}{(b^2-ac)(c^2-ab)},\quad O=-\frac{a(ab+bc+ca)(b^2+bc+c^2)}{(b^2-ac)(c^2-ab)}\]

 Note that 
 \[(ab+bc+ca)(b^2+bc+c^2)+(b^2-ac)(c^2-ab)=bc(a+b+c)^2\]The circumradius $R$ evaluates to
 \[R=\lvert O-a\rvert=\left\lvert \frac{bc(a+b+c)^2}{(b^2-ac)(c^2-ab)}\right\rvert\]And
 \[R^2=\frac{bc(a+b+c)^2}{(b^2-ac)(c^2-ab)}\cdot\frac{\frac{1}{bc}(\frac{1}{a}+\frac{1}{b}+\frac{1}{c})^2}{\left(\frac{1}{b^2}-\frac{1}{ac}\right)\left(\frac{1}{c^2}-\frac{1}{ab}\right)}\cdot\frac{a^2b^3c^3}{a^2b^3c^3}=\frac{bc(a+b+c)^2(ab+bc+ca)^2}{(b^2-ac)^2(c^2-ab)^2}\]Therefore, with\begin{gather*}O=\frac{-a(b^2+bc+c^2)(ab+bc+ca)}{(b^2-ac)(c^2-ab)},\quad I=\frac{bc}{2n+b+c},\\R=\frac{n(a+b+c)(ab+bc+ca)}{(b^2-ac)(c^2-ab)},\quad r=\frac{n+b+c}{2n+b+c},\end{gather*}it is enough to show that $\lvert O-I\rvert = \lvert R-r\rvert\iff (O-I)(\overline{O}-\overline{I})=(R-r)^2$.

 We make the following substitutions:\begin{gather*}S=b+c, \quad U=a+b+c,\quad V=ab+bc+ca,\\X=b^2+bc+c^2,\quad P=2n+b+c,\quad Q=(b^2-ac)(c^2-ab)\end{gather*}Substituting, we wish to prove that\begin{gather*}\left(-\frac{aXV}{Q}-\frac{n^2}{P}\right)\left(-\frac{UX}{Q}-\frac{1}{P}\right)=\left(\frac{nUV}{Q}-\frac{n+S}{P}\right)^2\iff\\(aXVP+n^2Q)(UXP+Q)=(nUVP-(n+S)Q)^2\iff\\UVP^2(aX^2-n^2UV)+PQ(X(n^2U+aV)+2n(n+S)UV)+(n^2-(n+S)^2)Q^2=0\tag{1}\end{gather*}To simplify, we express each variable in terms of $a,S,n$:\begin{gather*}U=a+S,\quad V=n^2+aS,\quad X=S^2-n^2\\P=2n+S,\quad Q=b^2c^2-a(b^3+c^3)+a^2bc=n^4-aS^3+3an^2S+a^2n^2\end{gather*}It follows that $n^2-(n+S)^2=-2nS-S^2=-PS$. Also,
 \[aX^2-n^2UV=a(S^4-2S^2n^2+n^4)-n^2(aS^2+a^2S+n^2S+an^2)=-S(n^4-aS^3+3an^2S+a^2n^2)=-SQ\]Substituting the two coefficients back to $(1)$ yields\begin{gather*}-SUVP^2Q+PQ(X(n^2U+aV)+2n(n+S)UV)-SPQ^2=0\iff\\X(n^2U+aV)+(2n^2+2nS-PS)UV-SQ=0\iff\\X(n^2U+aV)+(2n^2-S^2)UV-SQ=0\end{gather*}Expanding, we have\begin{align*}X(n^2U+aV)&=2an^2S^2+n^2S^3+a^2S^3-2an^4-n^4S-a^2n^2S\\(2n^2-S^2)UV&=2an^2S^2+2n^4S+2a^2n^2S+2an^4-aS^4-a^2S^3-n^2S^3-an^2S^2\\-SQ&=-n^4S+aS^4-3an^2S^2-a^2n^2S\end{align*}
Those terms cancel when added, as desired.

 


	\item (\textit{USAMO 2026}) A positive integer $n$ is called solitary if, for any nonnegative integers $a$ and $b$ such that $a+b=n$, either $a$ or $b$ contains the digit $1$ when written in base $10$. Determine, with proof, the number of solitary integers less than $10^{2026}$.


 



\textbf{Solution 1}

\textbf{Lemma 1: $n=2\cdot 10^k-1$ is solitary for all $k\in \mathbb{Z}_{\geq 0}$.}

 Suppose $a+b=n$. WLOG, $a\geq b$. Then, $2a\geq a+b=n\implies a \geq \frac{n}{2}$. Since $a, n$ are integers, and $n$ is odd, this implies $a \geq \frac{n+1}{2}=10^k$. But, $a\leq n = 2\cdot 10^k-1$, so $10^k \leq a \leq 2\cdot 10^k-1$, so the leading digit of $a$ is $1$. Thus, $n$ is solitary. $\square$

 
\textbf{Lemma 2: $m=2\cdot 10^k+n, n < 10^k$ is solitary if and only if $n$ is solitary.}

 First, suppose $m$ is solitary and $n$ is not. We write $n=a+b$ for some $a, b$ with no $1$'s in their decimal representation. Then, let $a' = a +2\cdot 10^k$ and $b' = b$. It follows that 
 \[a'+b' = a+2\cdot 10^k+b=2\cdot 10^k+n=m\]But $a'$ and $b'$ have no $1$'s in their decimal representation, since $10^k > n \geq a, b$, a contradiction. This proves the only if part.

 
Now, suppose $n$ is solitary, and $m$ is not. Write $m=a+b$ for some $a, b$ with no $1$'s in their decimal representation.

 Write $a = x\cdot 10^k+a', a'< 10^k$ (the leading digit of $a$ is $x$). Similarly, $b = y\cdot 10^k+b', b'< 10^k$. By definition, $a, b \neq 1$. But 
 \[x=y=0\implies a+b< 2\cdot 10^k < n \text{ and } x + y > 2\implies a+b\geq 3\cdot 10^k > m\] Thus, one of $x, y$ is $2$, and the other is $0$. WLOG $x=2, y=0$. Consider $a'$ and $b'$. We have 
 \[a'+b' = m- x\cdot 10^k -y\cdot 10^k =m - 2\cdot 10^k = n\] But $n$ is solitary, so either $a'$ or $b'$ has a $1$ in their decimal representation. But, that implies $a$ or $b$ also has a $1$, since $a', b' < 10^k$, a contradiction. This proves the if part. The lemma is complete. $\square$

 
\textbf{Lemma 3: For any integer $m \in [10^k, 10^{k+1}-1]$, if $m$ is not of the form $2\cdot 10^k - 1$ and not of the form $2\cdot 10^k + n$ (where $n < 10^k$), then $m$ is not solitary.}

 By Lemma 2, we can assume the leading digit of $m$ is not $2$. Now, take $a=10^k-1$ and $b = m-a$. 

 Let the decimal representations of $a, b$ be $a_{k-1}a_{k-2}\dots a_0$ and $b_kb_{k-1}\dots b_0$, respectively. It follows that $a_i=9$ for all $i$. 

 Also, $m\notin [2\cdot 10^k-1, 3\cdot 10^k-1]\implies b = m-a \notin [10^k, 2\cdot 10^k]$, so $b_k\neq 1$. 

 
Now, suppose $b_i = 1, i\in \{0, 1, \dots, k-1\}$. Increase $b_i$ to $2$ and lower $a_i$ to $8$. This increases $b$ by $10^i$ while decreasing $a$ by $10^i$, so the sum stays the same. Repeating this for all $i$ such that $b_i=1$, we have $a, b$ such that none of the digits are 1, $a+b=m$, so $m$ is not solitary. $\square$

 
Now, we use the lemmas to finish. Let $s_k$ be the number of solitary integers $< 10^k$, and let $S_k$ be the set of these integers, so $|S_k|=s_k$. By the lemmas 
 \[S_{k+1} = \{2\cdot 10^k-1\} \cup \{2\cdot 10^k+n : n\in S_k \}\cup S_k\] Since no elements are shared, we have $s_{k+1} = 1 + s_k+s_k=2s_k+1$. Also, $S_1 = \{1\}$, so $s_1 = 1$. It follows by induction $s_k = 2^k-1$. Thus, the answer is $\boxed{2^{2026}-1}$.

 



\textbf{Solution 2}

\textbf{Lemma 1:} $n$ must contain 1, otherwise it is not solitary.

 \textbf{Proof of Lemma 1:}

 Definition of solitary, choose $a = 0$ and $b = n$ and the result shows $n$ is not solitary.

 \textbf{Lemma 2:}: All such $n$ has the property all digits to the right of "1" are "9".

 \textbf{Proof of Lemma 2:}

 Consider vertical addition. Let $a_n$ and $b_n$ be the $n$th place value of such numbers $a$ and $b$. Then, $a_0 + b_0 = 9$. One can form such numbers using any orientation, either $a_n = 1$ and $b_n = 8$, vice versa, or some non-1 pair. Assume the worst case, where we construct sums of 9 that do not include 1. Then, if the leading digit itself is 1, then notice that one may "truncate" the 19 at the end of the string $1999\dots 9$. Because 9 is solitary, any single combination of $a_0$, $b_0$ result in one value being greater than 10, hence we must carry a value. Specifically, that implies the length of $a < b$ or vice versa. Hence, if the length $a < b$ is obtained, then notice that then one can let, at the end of the truncation, $b$ start with 9 and $a$ begin with 10, giving the 19 that was truncated at the beginning. This fulfills solitary properties.

 \textbf{Lemma 3:} All such $n$ has the property all digits to the left of "1" are either "0" or "2" or a combination of both.

 \textbf{Proof of Lemma 3:}

 We continue from where we left off in Lemma 2. We realize that $b$ must start with a 0 then move through with 9's. Afterwards, $a$ and $b$ are uniquely determined toward the left side of $a$'s 1 and $b$'s 0. Consider forming a sum of 11 instead of 1, and we prove that a number can be formed that is solitary. More importantly, if $a_n$ and $b_n$ do not contain 1 at all yet still sum to 11, then that 1 carries over. Now, we still fulfill the properties of \textbf{Lemma 2}, but the number is not solitary! But if we were to carry over that extra 1, then $a_{n+1}$ or $b_{n+1}$ can have either 1 or 9, from which we terminate, or continue analogously. However, when we end, one of $a$ or $b$ has 1, hence the number $n$ is solitary.

 \textbf{2026 USAMO/Q4:}

 Given the conditions, $n$ is \textit{solitary} if and only if

 
                \begin{itemize}
                    \item $n$ contains 1
                \end{itemize}


                
                \begin{itemize}
                    \item All numbers to the right of the 1 in $n$ are 9
                \end{itemize}


                
                \begin{itemize}
                    \item All numbers to the left of the 1 in $n$ are "2" or "0"
                \end{itemize}


                The last part is to count all such $n$. From any number, the digit 1 can be placed at an index $m$ from the left ($1 \leq m \leq 2026$). 

 There are exactly $2026 - m$ positions to the right of 1 that are all 9. Otherwise, there are exactly $m-1$ positions to the left of 1. Those have two choices each, giving $2^{m-1}$ cases. So, on any given integer with $m$ digits, we may form $2^{m-1}$ solitary integers.

 We sum over all $1 \leq m \leq 2026$, or

 
 \[\sum_{m = 1}^{2026} 2^{m-1} = 1+2+\dots+2^{2025} = \boxed{2^{2026}-1} \quad \square\].

 


	\item (\textit{USAMO 2026}) Let $ABC$ be a triangle. Points $D$, $E$, and $F$ lie on sides $BC$, $CA$, and $AB$, respectively, such that

 \[\angle AFE=\angle BDF=\angle CED.\]
Let $O_A$, $O_B$, and $O_C$ be the circumcenters of triangles $AFE$, $BDF$, and $CED$, respectively. Let $M$, $N$, and $O$ be the circumcenters of triangles $ABC$, $DEF$, and $O_AO_BO_C$, respectively. Prove that $OM=ON$.


 



\textbf{Solution 1}


                    \begin{figure}[H]
                        \centering
                        \includegraphics[width=0.5\textwidth]{Solution/USAMO_I_2026_5_Solution_Figure_1.png}
                    \end{figure}
                    

 
Let $P$ be the point where the circumcircles of $\triangle{AFE}, \triangle{BDF}, \triangle{CED}$ intersect, which we know exists because it is the Miquel point of the given configuration.

 \textbf{Claim 1:} We have $\triangle{ABC} \sim \triangle{EFD} \sim \triangle{O_AO_BO_C}$.

 \textbf{Proof:} For $\triangle{ABC} \sim \triangle{O_AO_BO_C}$, we have that $\angle{CAB} = \angle{EAP} + \angle{PAF} = \angle{O_CO_AP} + \angle{PO_AO_B} = \angle{O_CO_AO_B}$ and $\angle{BCA} = \angle{DCP} + \angle{PCE} = \angle{O_BO_CP} + \angle{PO_CO_A} = \angle{O_BO_CO_A}$ from $O_A$, $O_B$, $O_C$ being the centers of circumcircles $\triangle{AFE}, \triangle{BDF}, \triangle{CED}$.

 For $\triangle{ABC} \sim \triangle{EFD}$, we have from $\angle{AFE} = \angle{BDF} = \angle{CED}$ and $D, E, F$ lying on $BC, CA, AB$ respectively that 
 \[180^\circ = \angle{CED} + \angle{DEF} + \angle{FEA} = \angle{AFE} + \angle{CAB} + \angle{FEA} \implies \angle{DEF} = \angle{CAB}\] and analogously $\angle{FDE} = \angle{BCA}.$

 \textbf{Claim 2:} $P$ is the center of the spiral similarities $s_{1}, s_{2}$ sending $\triangle ABC$ to $\triangle EFD$ and $\triangle ABC$ to $\triangle O_AO_BO_C$ respectively.

 \textbf{Proof:} For $s_{1}$ note by cyclic quadrilaterals $AEPF, CDPE, BFPD$ that 
 \[\angle AFE = \angle BDF = \angle CED = \angle CPD = \angle BPF = \angle EPA,\] which we call $\theta$, and $\angle AEP = \angle BFP = \angle CDP$. Hence $\triangle AEP \sim \triangle BFP \sim \triangle CDP$ so $s_{1}$ rotates $\triangle ABC$ by $\theta$ and scales down by the common ratio $\frac{AP}{PE} = \frac{BP}{PF} = \frac{CP}{PD}$ to precisely yield $\triangle EFD$.

 For $s_{2}$ note that from $\angle PDC = \angle PEA = \angle PFB$ follows $\angle APO_A = \angle BPO_B = \angle CPO_C$ and $\angle PO_AA = \angle PO_BB = \angle PO_CC$, which means $\triangle APO_A \sim \triangle BPO_B \sim \triangle CPO_C$ and similarly suffices to imply that $P$ is the center of $s_{2}$.

 To finish, $s_{1}(M) = N$ and $s_{2}(M) = O$. Since $s_{1}(A) = E$ and $s_{2}(A) = O_A$ while $s_{1}, s_{2}$ share the same center, $\triangle MON \sim \triangle AO_AE$ from which $OM = ON$ because $O_AA = O_AE$.

 


	\item (\textit{USAMO 2026}) Let $a$ and $b$ be positive integers such that $\varphi(ab+1)$ divides $a^2+b^2+1$. Prove that $a$ and $b$ are Fibonacci numbers.


 Note: Euler’s totient function $\varphi(n)$ is the number of positive integers less than
or equal to $n$ that are relatively prime to $n$.


 



\textbf{Solution 1}

First, note that since $ab+1 > 2$, we have $2\mid \varphi (ab+1)$. Therefore, $2\mid a^2+b^2+1$, so one is even and the other is odd. 

 This implies $a^2+b^2+1\equiv 2\pmod{4}$, so $\nu _2 (\varphi(ab+1))=1$. To use this, note that $ab+1$ is odd.

 Let $ab+1 = p_1^{e_1}p_2^{e_2}\dots p_n^{e_n}$, where $p_1, p_2, \dots, p_n$ are odd primes. Then, 
 \[1 = \nu _2 (\varphi (ab+1)) = \nu _2 \left (\ \prod_{i=1}^{n} p_i^{e_i-1}(p_i-1) \right ) \geq n\] Thus, $ab+1 = p^k$ for some odd prime $p$ and some positive integer $k$. Furthermore, by the above reasoning, $p \equiv 3 \pmod{4}$.

 
\textbf{Case 1: $k\geq 2$}

 Since $p\nmid a$, we can write 
 \[ab+1 \equiv 0 \pmod{p^{k-1}}\implies b\equiv -a^{-1} \pmod{p^{k-1}}\] Now, $p^{k-1} \mid \varphi (ab+1) \mid a^2+b^2+1$, so we can substitute like so:
 \[a^2+b^2+1 \equiv 0 \pmod{p^{k-1}} \implies a^2+\frac{1}{a^2}+1\equiv 0 \pmod{p^{k-1}}\]
 \[\implies 4a^4+4a^2+4 \equiv 0 \pmod{p^{k-1}} \implies (2a^2+1)^2 \equiv -3 \pmod{p^{k-1}}\implies \left(\frac{-3}{p}\right) = 1\] where $\left(\frac{a}{p}\right)$ is the Legendre symbol. By quadratic reprocity, using the fact that $p\equiv 3 \pmod{4}$, 
 \[\left(\frac{-3}{p}\right) =\left(\frac{3}{p}\right) \left(\frac{-1}{p}\right) = -\left(\frac{3}{p}\right) = \left(\frac{p}{3}\right)\] Thus, $\left(\frac{p}{3}\right) = 1$, so $p\equiv 0 \pmod{3}\implies p=3$, or $p\equiv 1 \pmod{3}$. 

 
\textbf{Subcase 1: $p=3$}

 Note that $\left(\frac{-3}{9}\right) = -1$, so we must have $k-1\leq 1\implies k=2$. Thus, $ab+1 = 9$. It easy to check that $(a,b) = (1, 8), (8, 1)$ are the only solutions in this case. Since both $1$ and $8$ are Fibonacci numbers, this subcase is complete.

 
\textbf{Subcase 2: $p\equiv 1 \pmod{3}$} 

 Then, $3 \mid p-1 \mid \varphi (ab+1) \mid a^2+b^2+1$, so $3\nmid a, b$. However, $ab+1 \equiv p^k \equiv 1 \pmod{3}\implies 3 \mid ab$, a contradiction. Thus, this subcase yields no solutions.

 Putting all of the above together completes the first case.

 


 \textbf{Case 2: $k = 1$}

 Thus, $ab+1=p$, so $\varphi(ab+1)=ab$, and the problem condition is $ab \mid a^2+b^2+1$. Let $\frac{a^2+b^2+1}{ab}=t$, where $t$ is a positive integer. Let $(a,b)$ be the pair with the minimal value of $a+b$. WLOG, $a\geq b$. Then, 
 \[a^2-tab+b^2+1=0\] We know $a$ is one of the roots of this quadratic. Let $x$ be the other root of this quadratic. Then, $x = tb-a$ and $x = \frac{b^2+1}{a}$. The first equation proves that $x$ is an integer, and the second proves that $x$ is positive. Furthermore, if $b^2+1 <a^2$, then $x = \frac{b^2+1}{a} < a$, which gives $(x, b)$ as a valid pair, a contradiction. Thus, we must have $a=b$ for this minimal pair. However, $a^2 \mid 2a^2+1 \implies a=1$, so $(a, b) = (1, 1)$.

 This implies that $t=3$, since otherwise we could follow the above logic to obtain a contradiction by infinite descent.

 
We claim that all solutions are of the form $(a, b) = (F_i, F_{i-2})$ or $(F_{i-2}, F_i)$. Proceed by induction on $a+b$. Let $S_k$ be the set of solutions such that $a+b\leq k$. For the base case $k=2$, $S_k = \{(1, 1)\}$. Now, assume that all solutions in $S_{k-1}, k\geq 3$ are Fibonacci numbers. 

 Suppose we have a solution such that $a+b = k$. WLOG, $a\geq b$. By the above logic, we have a smaller $x$ such that $a = 3b - x$ and $x = \frac{b^2+1}{a}$. Since $a\geq 2$ and $x$ is an integer, we have $x = \frac{b^2+1}{a} < b$. Furthermore, since $(b, x)\in S_{k-1}$, we have $b = F_i$ and $x = F_{i-2}$. Thus, 
 \[a = 3b-x = 3F_i - F_{i-2} = 2F_i + F_{i-1} = F_{i+1}+F_i = F_{i+2}\] Therefore, $(a, b) = (F_{i+2}, F_{i})$. This completes the induction. 

 
Putting all this together finishes the second case. Finally, the proof is complete $\square$

 


\end{enumerate}